\documentclass[10pt,a4paper]{article}
\usepackage[margin=22mm]{geometry}
\usepackage[T1]{fontenc}
\usepackage{lmodern}
\usepackage{amsmath,amssymb,amsthm,booktabs,array,hyperref,enumitem}
\hypersetup{colorlinks=true,urlcolor=blue,linkcolor=blue,
  pdftitle={Exact Outer Multiset Dimension of the Three-by-Three Path Grid Strips},
  pdfauthor={Ryutaro Yonezu},
  pdfsubject={Computer-assisted mathematical preprint v1.0, 10 November 2026}}
\newcommand{\odim}{\operatorname{odim}}
\newcommand{\sq}{\mathbin{\square}}
\newcommand{\MS}{\mathcal M}
\newtheorem{theorem}{Theorem}
\newtheorem{lemma}{Lemma}
\newtheorem{proposition}{Proposition}
\title{\bfseries Exact Outer Multiset Dimension of the Three-by-Three Path Grid Strips}
\author{Ryutaro Yonezu\\\small Independent Researcher}
\date{10 November 2026 --- No. 62, Preprint v1.0}
\linespread{0.98}
\begin{document}
\maketitle
\begin{center}
\fbox{\begin{minipage}{0.93\linewidth}\centering\small
\textbf{COMPUTER-ASSISTED MATHEMATICAL PREPRINT --- v1.0}\\[2pt]
Independent AI differential audit through v0.3: author-reported PASS.\\
Windows PowerShell 5.1 reproduction of v0.3: author-supplied PASS.\\
Novelty: INCONCLUSIVE. Human peer review: NOT PERFORMED.\\
Copyright \copyright\ 2026 Ryutaro Yonezu. All rights reserved.
\end{minipage}}
\end{center}
\begin{abstract}
We give a computer-assisted exact classification of the outer multiset dimension of
$G_n=P_3\sq P_3\sq P_n$, for $n\ge2$. The value is $6$ for $n=3$,
$5$ for $n\in\{2,4,5,6,7,8\}$, and $4$ for $n\ge9$.
The lower bound $\odim(G_n)\ge4$ uses the separately published,
computer-assisted universal three-dimensional lower bound in the author's No.~59
preprint; it is not independently reproved here. The small exceptional cases
are established by exhaustive finite certificates, while the infinite upper
bound follows from four explicit landmarks, distance-sum identities, and two
bounded integer collision checks. The result is neither formally verified nor
human peer reviewed; novelty remains unconfirmed.
\end{abstract}

\section{Definition, scope, and relation to earlier work}
For a connected graph $G$ and $S\subseteq V(G)$, write
$\MS_S(v)=\{\!\{d_G(v,s):s\in S\}\!\}$ for the distance \emph{multiset},
with multiplicities retained and order forgotten. The set $S$ is
\emph{outer multiset resolving} if $\MS_S(u)\ne\MS_S(v)$ for all distinct
$u,v\in V(G)\setminus S$. Its minimum cardinality is $\odim(G)$.
The property is not generally monotone under adding landmarks; in particular,
absence of a $k$-landmark solution alone does not establish absence of every
smaller solution.

The two-path value $\odim(P_s\sq P_t)=3$ for $s\ge t\ge2$ is due to
Klav\v{z}ar--Kuziak--Yero~\cite{KKY2023}. Farhan--Klav\v{z}ar--Kuziak--Yero
list extension to more than two path factors as Problem 10 in their
2026 survey~\cite{FKKY2026}. This manuscript addresses only the subfamily
$P_3\sq P_3\sq P_n$, not general three-path products.

The author's No.~58~\cite{Yonezu58} examines $P_2\sq P_3\sq P_n$.
Our $G_2\cong P_2\sq P_3\sq P_3$ is the \emph{same graph} as the $n=3$
finite instance in the author's report on No.~58. The value $5$ is
also recorded in the short-grid table of No.~59~\cite{Yonezu59}.
The published theorem of No.~58 covers $n\ge6$. We include $G_2$ as
an overlap for completeness. No.~59 supplies the universal lower bound
\begin{equation}\label{eq:no59}
  \odim(P_a\sq P_b\sq P_c)\ge4\qquad (a,b,c\ge2).
\end{equation}
No.~59 is an author preprint with computer-assisted verification, not an
independently formalized theorem of the external literature. Our reliance
on it is explicit.

Identify $V(G_n)$ with $\{0,1,2\}^2\times\{0,\ldots,n-1\}$.
The graph distance is $d((x,y,z),(u,v,w))=|x-u|+|y-v|+|z-w|$.

\begin{theorem}\label{thm:main}
For every integer $n\ge2$,
\[
 \odim(P_3\sq P_3\sq P_n)=
 \begin{cases}
 6,&n=3,\\
 5,&n=2\text{ or }4\le n\le8,\\
 4,&n\ge9.
 \end{cases}
\]
The proof depends on the lower-bound result \eqref{eq:no59}, finite
exhaustive computation, and the exact finite certificates detailed below.
\end{theorem}

\section{Small parameters and finite certificates}
For each $n=2,\ldots,8$, the frozen C++17 verifier exhaustively tests
all $\binom{9n}{4}$ four-landmark sets and finds none outer multiset
resolving. For $n=3$, it also tests all $\binom{27}{5}$ five-landmark
sets, again finding none. This verifier represents a $k$-distance multiset
by its first $k$ integer power sums. Newton's identities imply that
these determine the multiset exactly; no probabilistic hash collision is
accepted as equivalence because unordered-set equality compares the full
power-sum array (the hash is only an index).

\begin{center}\small
\begin{tabular}{@{}rrrrrrrr@{}}\toprule
$n$&2&3&4&5&6&7&8\\\midrule
$\binom{9n}{4}$&3060&17550&58905&148995&316251&595665&1028790\\
\bottomrule\end{tabular}
\end{center}
The additional $n=3$ five-set total is $\binom{27}{5}=80730$.
Explicit upper witnesses are recorded verbatim in
\texttt{frozen/verify\_finite.cpp}, for sizes $5$ at $n=2,4,\ldots,8$
and size $6$ at $n=3$. Their distance multisets are checked directly.
With \eqref{eq:no59} excluding every size at most three, these computations
settle the minimum values for $n\le8$. The source and expected outputs are
frozen and accompanied by SHA256 digests.

\section{Four-landmark constructions for all \texorpdfstring{$n\ge9$}{n >= 9}}
For odd $n\ge9$ let $c=3$; for even $n\ge12$ let $c=6$. Set
\begin{equation}\label{eq:landmarks}
S_n=\{(0,0,0),(0,0,1),(0,0,n-1),(0,2,c)\}.
\end{equation}
At $n=10$ use instead the exceptional witness
\begin{equation}\label{eq:n10}
 S_{10}=\{(0,0,0),(0,0,6),(0,2,8),(2,2,5)\}.
\end{equation}
For $9\le n\le31$, the frozen Python certificate directly compares
all sorted four-distance signatures outside $S_n$ and finds them distinct.
For $n\ge32$, the following lemmas reduce infinitely many possible
collisions to small, explicit integer checks.

Let $p=x+y$, $a=2x+y$, and write the distances to the four landmarks
of \eqref{eq:landmarks} in the listed order as $A,B,C,D$. Then
\begin{equation}\label{eq:distances}
 (A,B,C,D)=(p+z,\ p+|z-1|,\ p+n-1-z,\ x+2-y+|z-c|).
\end{equation}
Their sum $T=A+B+C+D$ satisfies the exact identity
\begin{equation}\label{eq:sum}
 T=\begin{cases}
 n+c+2+2a,&z=0,\\
 n+c+2a,&1\le z<c,\\
 n-c+2(z+a),&z\ge c.
 \end{cases}
\end{equation}
Here $0\le a\le6$.

\begin{lemma}[Near-prefix reduction]\label{lem:prefix}
Suppose $n\ge32$ and two distinct vertices outside $S_n$ have equal
four-distance multisets. If either vertex has $z<c$, both have
$z\le c+7$.
On this prefix, $C$ is strictly larger than each of $A,B,D$.
Consequently any collision in this case is detected by equality of
\[
 \left(p-z,\;\operatorname{sort}\bigl(A,B,D\bigr)\right).
\]
The exact finite prefix check finds no such collision.
\end{lemma}
\begin{proof}
By \eqref{eq:sum}, a vertex with $z<c$ has $T\le n+c+14$.
For a vertex at $z'\ge c$, equality of the sums implies
$n-c+2z'\le n+c+14$, whence $z'\le c+7$.
If both vertices have $z<c$, the same bound holds trivially.
On $0\le z\le c+7\le13$ and $n\ge32$ we have
$C-A=n-1-2z>0$, and $C>B$ by inspection of $z=0$ or $z\ge1$.
Moreover
$C-D=n-3+2y-z-|z-c|>0$: for $z\le c$, it is at least
$n-3-c$, and for $z\ge c$, at least $n-c-17$.
Thus $C$ is the unique maximum. Equal multisets must have $C=C'$,
so $p-z=p'-z'$; the other three distances must agree as multisets.
The frozen exact checker enumerates every vertex in this prefix except
$(0,0,0),(0,0,1),(0,2,c)$ (the fourth landmark $(0,0,n-1)$ is
outside the prefix). It checks $96$ vertices for $c=3$, and $123$ for
$c=6$, and finds all indicated signatures unique.
\end{proof}

\begin{lemma}[Tail collision certificate]\label{lem:tail}
For $n\ge32$, no two distinct vertices with $z,z'\ge c$ have equal
four-distance multisets with respect to \eqref{eq:landmarks}.
\end{lemma}
\begin{proof}
Set $dx=x-x'$, $dy=y-y'$, and $t=p+z$.
By the third case of \eqref{eq:sum}, equality of sums implies
\begin{equation}\label{eq:dz}
 z'-z=2dx+dy.
\end{equation}
Subtract the common integer $t$ from each element of the two multisets.
With $h=n-1-2z$ this yields the necessary equality
\begin{equation}\label{eq:tail}
 \{\!\{0,-1,2-2y-c,h\}\!\}
 =\{\!\{dx,dx-1,dx+2-2y'-c,h-3dx-2dy\}\!\}.
\end{equation}
Since $dx,dy\in[-2,2]$, the first three elements of the left
multiset lie in $[-8,0]$ and those on the right in $[-10,2]$.
Also $|3dx+2dy|\le10$.
If $|h|>30$, the $h$-term is the unique extreme (largest if $h>30$,
smallest if $h<-30$), on both sides, forcing
$h=h-3dx-2dy$. The sole solution of $3dx+2dy=0$
with $dx,dy\in[-2,2]$ is $dx=dy=0$; then \eqref{eq:dz}
forces $z=z'$, contradicting distinctness.

It remains to consider $|h|\le30$. Since $h\equiv n-1\pmod2$,
there are 31 allowed even $h$ values when $n$ is odd ($c=3$), and
30 allowed odd $h$ values when $n$ is even ($c=6$).
For each $c$, exhaust the 72 ordered \emph{distinct} pairs
$(x,y)\ne(x',y')$ in $\{0,1,2\}^2$, together with every such $h$.
Sorting the four entries on both sides of \eqref{eq:tail}
produces no equality among $72\cdot31=2232$ checks for $c=3$,
respectively $72\cdot30=2160$ checks for $c=6$.
This is a safe overapproximation: the checker does \emph{not}
require the corresponding $z,z'$ to be feasible grid coordinates.
No floating-point arithmetic or random sampling is used.
\end{proof}

\begin{proof}[Proof of Theorem~\ref{thm:main}]
Equation \eqref{eq:no59} excludes all outer multiset resolving sets of size
at most three. The finite exhaustive computation and the explicit
upper witnesses settle $2\le n\le8$. For $9\le n\le31$, the direct
four-point certificate supplies an upper bound and \eqref{eq:no59} a
matching lower bound. For $n\ge32$, every pair of distinct exterior
vertices is either covered by Lemma~\ref{lem:prefix} or has both
$z,z'\ge c$ and is covered by Lemma~\ref{lem:tail}.
Thus \eqref{eq:landmarks} gives the upper bound four, which is sharp
by \eqref{eq:no59}.
\end{proof}

\section{Reproducibility and limitations}
The original exploration sources \texttt{verify\_finite.cpp} (C++17)
and \texttt{verify\_infinite.py} (Python standard library), together with
their two expected stdout records, are preserved byte for byte from
v0.3 in the accompanying v1.0 reproducibility package. The first
exhaustively enumerates the finite lower-bound cases, using exact integer
power sums. The second checks finite instances and the two analytical
collision reductions. The runner verifies the file inventory and SHA256
digests, compiles in a temporary directory, and compares the outputs with
the frozen records. Release preparation reproduced both outputs on Linux.

The author's supplied reports record a scope-limited independent AI
differential audit PASS through v0.3. The reported scope covers the model,
finite and infinite claims, editorial changes, package integrity, and privacy.
Earlier reports describe independent BFS checks for finite cases and a
rederivation of the infinite collision reduction. The v0.3 reviewer's
independent lower-bound tests were finite. The all-$n$ lower bound in this
paper relies on No.~59. The raw reviewer materials are private.

The author's supplied transcript records a Windows PowerShell 5.1
reproduction PASS for the exact v0.3 package on 10 November 2026.
It covers the inventory and SHA256 checks, the four frozen files, and
both computational certificates. The reproducibility package includes a
sanitized summary of that execution. The Windows result applies to v0.3.
We checked the v1.0 release package on Linux.

The literature search is incomplete and the novelty status is INCONCLUSIVE.
Human peer review and full formal verification have not been performed.

\begin{thebibliography}{9}\small
\bibitem{KKY2023}
S. Klav\v{z}ar, D. Kuziak, I. G. Yero,
``Further Contributions on the Outer Multiset Dimension of Graphs,''
\emph{Results in Mathematics} \textbf{78} (2023), Article 50.
\url{https://doi.org/10.1007/s00025-022-01829-8}.
\bibitem{FKKY2026}
M. Farhan, S. Klav\v{z}ar, D. Kuziak, I. G. Yero,
``Multiset resolvability parameters in graphs: A survey with new results
and open problems,'' arXiv:2607.10311v1 (2026), Problem 10.
\url{https://arxiv.org/abs/2607.10311}.
\bibitem{Yonezu58}
R. Yonezu, ``Exact Outer Multiset Dimension of the Three-Dimensional
Grid Family $P_2\sq P_3\sq P_n$,'' author preprint No.~58, v1.0 (2026).
\href{https://github.com/yonezaemon1-hub/no58-outer-multiset-dimension-3d-grid/releases/tag/v1.0}{GitHub release v1.0}.
\bibitem{Yonezu59}
R. Yonezu, ``Outer Multiset Dimension of Three-Dimensional Grids:
A Universal Lower Bound and an Exact Two-Parameter Classification,''
author preprint No.~59, v1.0 (2026), universal lower-bound theorem.
\href{https://github.com/yonezaemon1-hub/no59-outer-multiset-dimension-3d-grid-classification/releases/tag/v1.0}{GitHub release v1.0}.
\end{thebibliography}
\end{document}
